% LayoutTranslateBench --- arXiv preprint (cs.CL)
% arXiv-safe: standard `article` class, no conference style files.
% Build: pdflatex main && bibtex main && pdflatex main && pdflatex main
\documentclass[11pt]{article}

\usepackage[utf8]{inputenc}
\usepackage[T1]{fontenc}
\usepackage{amsmath}
\usepackage{amssymb}
\usepackage{booktabs}
\usepackage{graphicx}
\usepackage{array}
\usepackage{multirow}
\usepackage{xcolor}
\usepackage[margin=1in]{geometry}
\usepackage[numbers,sort&compress]{natbib}
\usepackage{hyperref}
\hypersetup{
  colorlinks=true,
  linkcolor=blue!55!black,
  citecolor=blue!55!black,
  urlcolor=blue!55!black,
}

% Compact column type for numeric table cells.
\newcolumntype{R}{>{\raggedleft\arraybackslash}p{1.6cm}}

\title{LayoutTranslateBench: A Benchmark for Document\\ Translation with Layout Preservation}

\author{Laurent Scotto}
\date{Dataset v0.1.7 \quad / \quad Spec v0.1 \\[2pt] \today}

\begin{document}
\maketitle

\begin{abstract}
Machine translation is evaluated almost exclusively on plain text, yet most
documents that people translate in practice---scanned forms, certificates,
contracts, scientific papers, news pages---carry information in their
\emph{layout}: the spatial arrangement of text regions and the order in which
those regions are meant to be read. Current document-translation tools
routinely degrade this layout, and there is no public benchmark that scores
layout fidelity and reading order \emph{jointly} with translation quality, so
there is no shared target to optimize. We introduce \textbf{LayoutTranslateBench
(LTB)}, a public benchmark for document translation with layout preservation.
LTB pairs each source document with annotated text regions---polygons or
bounding boxes, a document-global reading-order index, and reference
translations from certified translators or FLORES-200 parallel data---and asks
a system to declare, for every region, both the translated text and the box
where it placed that text. From this dual reporting LTB computes three metrics:
character-level \textbf{chrF} with a language-detection gate, \textbf{layout
IoU}, and a coverage-aware reading-order \textbf{Kendall~$\tau$}; it combines
them into a single composite, \textbf{LTB-100}, and reports COMET-Kiwi-22 as an
optional reference-free second metric on a parallel leaderboard. Version 0.1.7
of the dataset contains 59 documents across 9 categories and 16 English-source
language pairs. We evaluate seven systems in two classes---oracle-layout
text-quality upper bounds (DeepL Text API, NLLB-200-distilled-600M,
Helsinki-NLP/opus-mt, and an identity baseline) and end-to-end systems
(identity, Qwen3-VL-2B-Instruct, and Florence-NLLB)---and report results on both
metrics with bootstrap confidence intervals. We find that the chrF--COMET gap is
largest on non-Latin script pairs, that one pair (en$\rightarrow$uz) is a
structural zero for 600M-class open models, and that a large gap separates
oracle-layout upper bounds from realistic end-to-end systems. We release the
code (Apache-2.0), the dataset (CC-BY-4.0 / CC-BY-SA-4.0), the leaderboard, and
a public methodology roadmap that tracks every known limitation.
\end{abstract}

\section{Introduction}
\label{sec:intro}

A large fraction of real-world translation demand is for \emph{documents}, not
free-standing sentences: immigration paperwork, certified legal translations,
localized scientific papers, corporate compliance material, and e-commerce
listings all need a target-language artifact that resembles the source. For
these inputs, the spatial layout is part of the message. A two-column scientific
paper read column-major, a government form whose checkboxes sit beside their
labels, a certificate whose seal and signature line have fixed positions---in
each case the arrangement of text regions and their reading order carry meaning
that a flat string does not.

Current tools handle this poorly. Plain-text translators (DeepL, Google
Translate text mode) translate strings well but discard layout entirely.
Document-mode tools (DeepL Documents, Google Translate documents, vision-language
chat assistants) accept a document but frequently degrade its layout:
text expands or contracts and overflows its box, multi-column reading order
collapses, fonts are flattened, and right-to-left targets are mirrored
inconsistently. These failures are visible to any user, yet they are
\emph{invisible to standard MT evaluation}, which scores text against a
reference and never inspects the page.

The missing ingredient is a benchmark. WMT shared tasks~\citep{kocmi2023wmt}
measure text-only translation quality; document-parsing benchmarks such as
OmniDocBench~\citep{ouyang2024omnidocbench} measure extraction from documents
but not translation; OCR benchmarks measure recognition but not layout-preserved
output in a target language. No public benchmark scores layout fidelity,
reading-order preservation, and text quality \emph{together}. Without such a
benchmark there is no single number for a layout-preserving translation system
to optimize, and no objective way to compare systems.

\paragraph{Contributions.} This paper introduces \textbf{LayoutTranslateBench
(LTB)} and makes the following contributions.
\begin{itemize}
  \item \textbf{A task formulation} for document translation with layout
  preservation: a system must declare, per source region, both the translated
  text and the bounding box where it placed that text. This dual reporting is
  what makes joint scoring of translation and layout possible
  (Section~\ref{sec:benchmark}).
  \item \textbf{A dataset} (v0.1.7): 59 annotated documents across 9 categories
  and 16 English-source language pairs, with text regions, polygons/bounding
  boxes, reading-order indices, and reference translations whose provenance and
  per-document license are fully recorded (Section~\ref{sec:benchmark}).
  \item \textbf{A metric suite}: chrF with a language-detection gate, layout
  IoU, and a coverage-aware reading-order Kendall~$\tau$, combined into the
  \textbf{LTB-100} composite, with COMET-Kiwi-22 as an optional reference-free
  second metric and bootstrap confidence intervals on every score
  (Section~\ref{sec:metrics}).
  \item \textbf{A reference evaluation} of seven systems in two classes
  (oracle-layout upper bounds and end-to-end systems), surfacing a chrF--COMET
  gap concentrated on non-Latin scripts, a structural-zero pair, and a large
  oracle-versus-end-to-end gap (Section~\ref{sec:experiments}).
  \item \textbf{A public methodology roadmap} that tracks every known
  limitation of the benchmark and its fix status, in the spirit of an honest,
  self-critical research artifact (Sections~\ref{sec:limitations}
  and~\ref{sec:conclusion}).
\end{itemize}

LTB is a research artifact. The code is released under Apache-2.0 and the
dataset under CC-BY-4.0 (author-curated and ML-curated documents) or
CC-BY-SA-4.0 (FLORES-derived documents), with per-document licenses recorded in
the dataset manifest.

\section{Related Work}
\label{sec:related}

\paragraph{Text-only MT evaluation.} The dominant paradigm evaluates
translation on plain text. Surface-overlap metrics---BLEU~\citep{papineni2002bleu}
and its standardized variant SacreBLEU~\citep{post2018sacrebleu}, and
chrF~\citep{popovic2015chrf,popovic2017chrfpp}---compare a hypothesis string to
a reference string. Learned metrics such as COMET~\citep{rei2020comet} and the
reference-free COMET-Kiwi~\citep{rei2022cometkiwi} correlate better with human
judgment. The WMT shared tasks~\citep{kocmi2023wmt} are the standard venue for
comparing systems under these metrics. None of this machinery inspects a
document page: layout fidelity and reading order are out of scope by
construction. LTB reuses chrF and COMET-Kiwi as its \emph{text} component but
adds layout and reading-order axes that text-only evaluation cannot express.

\paragraph{Document parsing and OCR benchmarks.}
OmniDocBench~\citep{ouyang2024omnidocbench} benchmarks PDF document parsing
across diverse layouts---it scores how well a system \emph{extracts} structured
content (Markdown/HTML) from a document. OCR benchmarks such as socOCRbench
\citep{sococrbench} score multi-region text recognition quality. Both measure a
source-language extraction step; neither scores a translation, nor whether a
translated document preserves the source layout. These benchmarks are
complementary to LTB: a system can extract a document well and still translate
it badly or destroy its layout when re-rendering.

\paragraph{Document reconstruction.}
OmniDoc-TokenBench~\citep{omnidoctokenbench} evaluates tokenizer/autoencoder
reconstruction fidelity on text-heavy documents---a generative-fidelity signal,
not an end-to-end translation evaluation. LTB instead scores the full
document-translation task.

\paragraph{Layout-aware models.} Document-understanding models such as
LayoutLM~\citep{xu2020layoutlm} jointly model text and 2-D layout, and
OCR-grounded vision models such as Florence-2~\citep{xiao2024florence2} produce
text with pixel-precise regions. These are \emph{systems} that could be
evaluated; they are not benchmarks for layout-preserving translation. LTB
provides the evaluation target such systems currently lack.

\paragraph{Positioning.} Table~\ref{tab:related} summarizes the gap. To our
knowledge LTB is the first public benchmark that scores translation quality,
layout fidelity, and reading-order preservation \emph{jointly}, on documents,
with a single composite score. We deliberately keep text, layout, and order as
\emph{separate} reported axes---rather than collapsing them into one per-pixel
bilingual signal as proposed for document-level visual metrics---so that a
system's failure mode is diagnosable from its score profile.

\begin{table}[t]
\centering
\small
\begin{tabular}{@{}p{3.0cm}p{5.0cm}p{5.0cm}@{}}
\toprule
\textbf{Benchmark} & \textbf{What it measures} & \textbf{What LTB adds} \\
\midrule
OmniDocBench & Document parsing / OCR --- extraction to Markdown/HTML, not
translation & Translation quality and layout scoring in a target language \\
\addlinespace
socOCRbench & Multi-region OCR recognition quality & Layout-preserved output,
not just source-language extraction \\
\addlinespace
WMT shared tasks & Text-only translation quality & Document layout and reading
order as first-class signals \\
\addlinespace
OmniDoc-TokenBench & Autoencoder reconstruction of text-heavy documents &
End-to-end translation evaluation \\
\bottomrule
\end{tabular}
\caption{LTB positioned against related benchmarks. The benchmarks are
complementary; LTB is the first to score translation, layout, and reading order
jointly.}
\label{tab:related}
\end{table}

\section{The LTB Benchmark}
\label{sec:benchmark}

\subsection{Task formulation}

A document in LTB is annotated as a set of \emph{text regions}. A region is a
polygon (or axis-aligned bounding box) enclosing a contiguous run of text in a
single logical block---a heading, a paragraph, a table cell, a stamp, a
signature line. Every ground-truth region carries: the source text; a reference
translation for each covered target pair; style hints (font-family class,
size hint, color, background); a 0-based document-global reading-order index;
and a layout class (e.g.\ \texttt{single-column}, \texttt{two-column},
\texttt{form-field}, \texttt{table-cell}, \texttt{caption}, \texttt{header},
\texttt{footer}).

A submission for one (system, language pair) is a JSON Lines file: one line per
document, each line listing predicted regions. For every region the system
declares (i)~the translated text it produced and (ii)~the bounding box where it
placed that text in its rendered output, plus a reading-order index. This
\emph{dual reporting}---text and box together---is the design choice that lets
LTB score translation quality and layout fidelity simultaneously from a single
submission file.

\subsection{Dataset (v0.1.7)}

Version 0.1.7 of the dataset comprises \textbf{59 documents} across 9
categories. Table~\ref{tab:categories} gives the actual category distribution.
The \texttt{ocr-document} category dominates at 39 of 59 documents (66\%); we
treat this imbalance as a known limitation (Section~\ref{sec:limitations}) and
target $\geq 15$ documents per category, all author-curated or
certified-translator sourced, for v0.2 (200 documents).

\begin{table}[t]
\centering
\small
\begin{tabular}{@{}lrl@{}}
\toprule
\textbf{Category} & \textbf{$n$ (v0.1.7)} & \textbf{Provenance} \\
\midrule
ocr-document       & 39 & ML-curated (multilingual-document corpus), script-validated \\
magazine-news      & 11 & 1 author-curated + 10 FLORES-200 (Wikinews) \\
scientific-paper   &  2 & author-curated \\
gov-form           &  2 & author-curated \\
certificate        &  1 & author-curated \\
receipt-invoice    &  1 & author-curated \\
business-letter    &  1 & author-curated \\
handwritten-mixed  &  1 & author-curated \\
legal-contract     &  1 & author-curated \\
\midrule
\textbf{Total}     & \textbf{59} & \\
\bottomrule
\end{tabular}
\caption{Document categories in LTB v0.1.7 and their provenance.
\texttt{ocr-document} is over-represented at 66\%; v0.2 rebalances to
$\geq 15$ documents per category.}
\label{tab:categories}
\end{table}

\paragraph{Language pairs.} LTB v0.1.7 covers \textbf{16 English-source language
pairs}, organized as a frozen \emph{core} set of 8 and an \emph{extension} set
of 8:
\begin{itemize}
  \item \textbf{Core 8} (frozen at spec v0.1, suitable for headline LTB-100
  reporting): en$\rightarrow$es, en$\rightarrow$de, en$\rightarrow$zh,
  en$\rightarrow$ar, en$\rightarrow$ja, en$\rightarrow$fr, en$\rightarrow$th,
  en$\rightarrow$ms. These were chosen to stress distinct layout effects: text
  expansion (German, $\sim$30\% longer), contraction (Chinese, $\sim$30--50\%
  shorter), right-to-left mirroring (Arabic), and script mixing (Japanese).
  \item \textbf{Extension 8} (added in v0.1.6, suitable for testing per-pair
  coverage of large multilingual systems but at sample sizes that do not
  support tight ranking): en$\rightarrow$ru, en$\rightarrow$ko,
  en$\rightarrow$vi, en$\rightarrow$id, en$\rightarrow$ur, en$\rightarrow$uz,
  en$\rightarrow$kk, en$\rightarrow$zh-tw.
\end{itemize}
At v0.1.7 every pair has at least 10 certified-translator-grade reference
documents. Per-pair sample sizes are $N=20$ for six core pairs (en-es, en-de,
en-ar, en-fr, en-th, en-ms), $N=27$ for en-zh, $N=28$ for en-ja, $N=10$ for
en-ru (FLORES-200 only), and $N=13$ for each of the remaining seven extension
pairs (10 FLORES-derived + 3 ML-curated). Confidence intervals on the extension
pairs are deliberately wide---they are coverage proofs, not benchmark-quality
samples.

\paragraph{Provenance and licensing.} Reference quality is mixed by document and
is fully transparent in each annotation's \texttt{provenance.grade} field.
Author-curated references are single translations at a native-speaker,
non-professional level; ML-curated references come from a multilingual-document
corpus; certified-translator references are derived from the FLORES-200
parallel evaluation data~\citep{flores200}. The code is Apache-2.0; the dataset is CC-BY-4.0 for
author-curated and ML-curated documents and CC-BY-SA-4.0 for FLORES-derived
documents, with per-document licenses recorded in \texttt{data/manifest.json}.

\paragraph{Reference-script validation.} While investigating an anomaly on
en-uz / en-ru, spot-checks revealed that the mined ML-curated source corpus
contained region-level reference entries labeled as one language but written in
a different script entirely (e.g.\ Simplified Chinese text under an en-ru
label). A validation pass (\texttt{scripts/validate\_extension\_refs.py})
detects script-mismatched references region-by-region and drops them; the
document itself is retained as layout ground truth, and the coverage-aware
scorer (Section~\ref{sec:metrics}) prevents the dropped regions from inflating
the affected pair's average. This pass removed 103 region references across the
ML-curated subset and reduced en-ru to zero documents at v0.1.6.4; v0.1.7
restored en-ru to $N=10$ with FLORES-200 references. We recommend any external
corpus adopted by LTB be run through this validation pass at ingest.

\paragraph{Held-out split.} To discourage benchmark overfitting, 20\% of
documents per category are held out and never published in raw form. Submitters
score the public split locally; maintainers re-score the held-out split before
promoting a result to the leaderboard. The split rotates on a quarterly cadence.

\subsection{Annotation schema}

Each ground-truth document is a JSON record of regions. Region geometry is a
polygon or an axis-aligned bounding box in pixel coordinates relative to the
rendered source page. The reading-order index is 0-based and document-global, so
multi-column and footnote structure is expressible. Reference translations are
attached per region per language pair. Style hints are coarse (font-family
\emph{class} rather than a specific typeface) because the v0.1 metrics do not
score typography; the v0.2 visual-fidelity metric will make finer use of them.

\section{Metrics}
\label{sec:metrics}

LTB scores a system on three primary axes---text quality, layout fidelity, and
reading-order preservation---and combines them into the LTB-100 composite. For
each (system, language pair, document) triple the system produces predicted
regions; the scorer matches predicted regions to ground-truth regions, scores
each match, aggregates per document, then averages across documents and language
pairs.

\subsection{Region matching}

A predicted region is matched to a ground-truth region in two phases:
(i)~an exact \texttt{region\_id} match, honored only when the two boxes also
overlap (IoU~$\geq 0.10$), so that an end-to-end system whose region ids happen
to collide with the ground-truth namespace is not mis-paired to a
non-overlapping region; then (ii)~greedy IoU matching for any remaining regions,
with the same minimum IoU of $0.10$, candidate pairs sorted by descending IoU
and matched without replacement. Predicted regions that fail to match contribute nothing positive.
Ground-truth regions with no match contribute zero to chrF and IoU and are
excluded from the reading-order sequence.

\subsection{Text quality: chrF with a language-detection gate}

The text metric is chrF$_2$---character-level F-score with $\beta=2$ over
character $n$-grams for $n \in \{1,\dots,6\}$, the WMT-standard variant
\citep{popovic2015chrf}. For a ground-truth region with reference translation
$r$ and matched predicted text $h$,
\begin{equation}
\mathrm{chrF}_2(h, r) \;=\; \frac{1}{6}\sum_{n=1}^{6} F_2\big(\text{$n$-grams}(h),\,\text{$n$-grams}(r)\big),
\qquad
F_2(P,R) = \frac{5\,P\,R}{4\,P + R}.
\end{equation}
Per-document chrF is the \emph{area-weighted} mean of region chrF scores, so a
large body paragraph counts for more than a small stamp.

\paragraph{Language-detection gate.} chrF rewards character $n$-gram overlap
\emph{regardless of whether the output is in the target language}. In LTB v0.1
this was an artifact: an identity baseline that returns English text scored
chrF~$\approx 25$ on Latin-script targets via shared Latin characters, copied
proper nouns, and verbatim numbers and dates. Starting in v0.1.1, before
computing chrF for a region the predicted text is checked against the target
language; if it is confidently \emph{not} in the target language, chrF is forced
to $0$ for that region. The gate has two layers:
\begin{enumerate}
  \item A \textbf{script-based pre-filter} for non-Latin targets (zh, ja, ar,
  th, and the non-Latin extension scripts): the fraction of alphabetic
  characters falling in the target's Unicode script block is computed; $\geq
  50\%$ accepts, $<50\%$ rejects. This is more robust than a statistical
  detector on short strings.
  \item \textbf{Statistical language detection} for Latin-script targets (es,
  de, fr, ms, and Latin-script extension pairs), where script alone is
  insufficient.
\end{enumerate}
The gate is deliberately \emph{permissive}: short strings ($<4$ alphabetic
characters), numbers, currency, and pure punctuation pass unchecked; detection
failures fail open; and for en-ms both Malay and Indonesian are accepted because
they are mutually intelligible and not reliably separable on short strings. The
gate therefore fires only on \emph{confidently wrong-language} predictions and
does not penalize paraphrase variation or technical terms that legitimately
remain in the source language.

\subsection{Layout fidelity: IoU}

For each matched region pair $(g, p)$ the scorer computes axis-aligned
bounding-box intersection-over-union,
\begin{equation}
\mathrm{IoU}(g, p) \;=\; \frac{\mathrm{area}(g \cap p)}{\mathrm{area}(g \cup p)}.
\end{equation}
Per-document layout IoU is the area-weighted mean over all ground-truth regions;
unmatched ground-truth regions count as $\mathrm{IoU}=0$. This rewards systems
that place translated text in the same regions of the page as the source.

\subsection{Reading order: coverage-aware Kendall $\tau$}

For each document, after region matching, the scorer extracts two parallel
sequences---the ground-truth reading orders of the matched regions in document
order, and the predicted reading orders of those same regions---and computes
Kendall's $\tau$-b~\citep{kendall1938tau} between them ($\tau$-b handles ties),
normalized to $[0,1]$ via $\tau_{\mathrm{norm}} = (\tau + 1)/2$.

\paragraph{Coverage-aware correction.} A naive $\tau$ has a partial-coverage
hole: a system that returns only one region out of seven trivially scores
$\tau_{\mathrm{norm}}=1.0$, since a single-element sequence is in order by
definition. To remove this free credit, LTB scales $\tau$ by coverage:
\begin{equation}
\tau_{\mathrm{final}} \;=\; \tau_{\mathrm{norm}} \times
\min\!\left(1.0,\; \frac{n_{\mathrm{matched}}}{n_{\mathrm{gt}}}\right),
\end{equation}
where $n_{\mathrm{matched}}$ is the number of predicted regions that matched a
ground-truth region and $n_{\mathrm{gt}}$ is the document's total ground-truth
region count. Under this rule the one-of-seven fallback scores
$1.0 \times (1/7) \approx 0.14$ rather than $1.0$.

\subsection{Aggregation and the LTB-100 composite}

Per-region scores are combined by area-weighted mean into a per-document score;
per-document scores are combined by arithmetic mean into a per-language-pair
score; per-language-pair scores are combined by arithmetic mean (over pairs with
$n_{\mathrm{docs}} > 0$) into the overall score. Each language pair is weighted
\emph{equally} in the overall score regardless of document count, which protects
the benchmark against drift if documents are added unevenly across pairs.

The composite, \textbf{LTB-100}, is a weighted linear combination of the three
normalized metrics. The v0.1.x formula in production is
\begin{equation}
\mathrm{LTB\text{-}100}_{\,\mathrm{v0.1.x}}
= 100 \times \big(0.50 \cdot \mathrm{chrF}/100
+ 0.30 \cdot \mathrm{IoU}
+ 0.20 \cdot \tau \big).
\label{eq:ltb100v01}
\end{equation}
The weighting is intentionally biased toward text quality (50\%) because
translation correctness remains the dominant signal at this stage of the field.
The v0.2 formula adds an OCR round-trip term and an LPIPS visual-fidelity term
(Section~\ref{sec:conclusion}) and rebalances the existing weights:
\begin{equation}
\mathrm{LTB\text{-}100}_{\,\mathrm{v0.2}}
= 100 \times \big(0.40 \cdot \mathrm{chrF}/100
+ 0.25 \cdot \mathrm{IoU}
+ 0.15 \cdot \tau
+ 0.10 \cdot \mathrm{OCR}
+ 0.10 \cdot \mathrm{LPIPS} \big).
\label{eq:ltb100v02}
\end{equation}
Submitted v0.1 results remain valid under v0.2 and are displayed with an
explicit version label.

\paragraph{Weight choice.} The 50/30/20 weighting is not arbitrary: it was
validated post-hoc against three alternatives---balanced (40/40/20),
text-heavy (60/20/20), and uniform (33/33/33). System rankings are
\emph{identical} under all four schemes (rank correlation
$\tau_{\mathrm{norm}}=1.0$ between every pair of rankings) on the systems
evaluated, so weight choice does not change leaderboard ordering; the 50/30/20
weights are retained because they communicate the intended emphasis. If a future
system enters with chrF/IoU/$\tau$ in markedly different relative magnitudes,
the ablation is to be re-run, and if rankings diverge the weights must be
validated against human judgment.

\subsection{Bootstrap confidence intervals}

Every LTB-100 score is reported with a 95\% bootstrap percentile confidence
interval~\citep{efron1979bootstrap}: 1000 resamples with replacement from the
per-document scores, with a fixed seed (42) for reproducibility. The leaderboard
displays scores as \texttt{point [CI low, CI high]}. Resampling over documents
rather than reporting a bare point estimate follows standard practice for
machine-translation evaluation, where small test sets make point differences
unreliable~\citep{koehn2004statistical}.

\subsection{COMET-Kiwi-22 as an optional second metric}

chrF has well-known weaknesses---it is paraphrase-blind and adequacy-blind. LTB
therefore supports \textbf{COMET-Kiwi-22}~\citep{rei2022cometkiwi}, a
reference-free neural quality-estimation metric, as an optional second text
metric (\texttt{ltbench score --text-metric comet-kiwi}). COMET-Kiwi produces a
score in $[0,100]$ that substitutes for chrF in the composite formula, yielding
a \emph{parallel} leaderboard. chrF remains the default because it is fast,
deterministic, and CPU-only, whereas COMET-Kiwi requires a separate Python
environment and a gated model download. Reporting both metrics is itself a
design choice: rankings can differ between them, and where they differ is
informative (Section~\ref{sec:experiments}).

\subsection{Oracle-layout versus end-to-end systems}

A runner declares its \texttt{system\_type}. An \textbf{end-to-end} runner
produces its own bounding boxes from the source image---this is the realistic
measurement. An \textbf{oracle-layout} runner copies ground-truth bounding boxes
as its predictions and translates only the text---this is a text-quality
\emph{upper bound}, not a real-world measurement of the underlying product. The
leaderboard segregates the two classes into separate tables. Mixing them
misleads readers: an oracle-layout score for a commercial translator measures
``that translator's text quality assuming a perfect layout extractor,'' a
capability the product itself does not provide. We stress that oracle-layout and
end-to-end runners receive \emph{symmetric per-region inputs}---both are called
per region---so comparisons \emph{within} a class are fair.

\section{Experiments and Results}
\label{sec:experiments}

\subsection{Evaluated systems}

We evaluate seven systems in two classes.

\paragraph{Oracle-layout (text-quality upper bounds).}
\begin{itemize}
  \item \textbf{DeepL Text API} --- a commercial MT system, called per region
  via its Text API. Covers 6 of 16 pairs (no Thai, no Malay, plus uncovered
  extension pairs).
  \item \textbf{NLLB-200-distilled-600M}~\citep{nllb2022} --- an open-weight
  multilingual model (CC-BY-NC-4.0, research-only) covering all 16 pairs.
  \item \textbf{Helsinki-NLP/opus-mt}~\citep{tiedemann2020opusmt} --- open
  Marian-family models~\citep{junczys2018marian}, commercial-safe
  (Apache-2.0 / CC-BY-4.0), covering 16 of 16 pairs through a mix of per-pair
  and multi-target router models.
  \item \textbf{Identity baseline} --- returns the source text unchanged in the
  source boxes; a trivial lower bound on text quality and an upper bound on
  layout.
\end{itemize}

\paragraph{End-to-end.}
\begin{itemize}
  \item \textbf{Identity baseline} --- as above; under end-to-end accounting it
  also receives the source boxes, so it scores the layout ceiling with
  near-zero text quality.
  \item \textbf{Qwen3-VL-2B-Instruct}~\citep{qwen3vl2025} --- a zero-shot
  vision-language model that must extract its own regions and translate them.
  \item \textbf{Florence-NLLB} --- a two-stage end-to-end pipeline:
  Florence-2~\citep{xiao2024florence2} grounded OCR for region detection,
  followed by NLLB-200~\citep{nllb2022} for per-region translation. Described
  in Section~\ref{sec:florence}.
\end{itemize}

\subsection{Leaderboard A: chrF}

Table~\ref{tab:chrf-e2e} reports end-to-end systems and Table~\ref{tab:chrf-oracle}
reports oracle-layout upper bounds, scored with chrF. LTB-100 cells show the
point estimate with a 95\% bootstrap confidence interval.

\begin{table}[t]
\centering
\small
\begin{tabular}{@{}llRRRc@{}}
\toprule
\textbf{System} & \textbf{LTB-100 [95\% CI]} & \textbf{chrF} &
\textbf{IoU} & \textbf{$\tau$} & \textbf{Coverage} \\
\midrule
identity-baseline v0.1.0      & 50.38 [50.2, 50.6] & 0.62 & 1.0000 & 1.0000 & 16/16 \\
florence-nllb v0.1.0          & 39.79 [36.6, 42.9] & 37.93 & 0.1541 & 0.8104 & 8/16 \\
qwen3-vl-2b-instruct v0.1.0   & 14.82 [12.1, 17.4] & 5.13 & 0.1084 & 0.4504 & 8/16 \\
\bottomrule
\end{tabular}
\caption{Leaderboard A (chrF) --- end-to-end systems. These runners produce
their own bounding boxes; this is the realistic real-world score.
Florence-NLLB is detailed in Section~\ref{sec:florence}.}
\label{tab:chrf-e2e}
\end{table}

\begin{table}[t]
\centering
\small
\begin{tabular}{@{}llRRRc@{}}
\toprule
\textbf{System} & \textbf{LTB-100 [95\% CI]} & \textbf{chrF} &
\textbf{IoU} & \textbf{$\tau$} & \textbf{Coverage} \\
\midrule
deepl-text-oracle v0.1.0   & 78.20 [75.2, 81.7] & 56.40 & 1.0000 & 1.0000 & 6/16 \\
nllb-text-oracle v0.1.0    & 73.71 [72.5, 74.9] & 47.61 & 1.0000 & 1.0000 & 16/16 \\
opus-mt-text-oracle v0.1.0 & 68.60 [67.1, 70.0] & 37.26 & 1.0000 & 1.0000 & 16/16 \\
\bottomrule
\end{tabular}
\caption{Leaderboard A (chrF) --- oracle-layout reference. Each runner is given
ground-truth bounding boxes as predictions and translates only the text. These
are \emph{upper bounds} on text quality, not realistic end-to-end measurements.}
\label{tab:chrf-oracle}
\end{table}

\subsection{Leaderboard B: COMET-Kiwi-22}

Tables~\ref{tab:comet-e2e} and~\ref{tab:comet-oracle} report the same systems
under the reference-free COMET-Kiwi-22 metric. COMET-Kiwi replaces chrF in the
composite; IoU and $\tau$ are unchanged.

\begin{table}[t]
\centering
\small
\begin{tabular}{@{}llRRRc@{}}
\toprule
\textbf{System} & \textbf{LTB-100 [95\% CI]} & \textbf{COMET-Kiwi} &
\textbf{IoU} & \textbf{$\tau$} & \textbf{Coverage} \\
\midrule
identity-baseline v0.1.0    & 50.47 [50.3, 50.7] & 0.81 & 1.0000 & 1.0000 & 16/16 \\
florence-nllb v0.1.0        & 48.11 [44.1, 51.7] & 54.56 & 0.1541 & 0.8104 & 8/16 \\
qwen3-vl-2b-instruct v0.1.0 & 18.43 [14.6, 22.3] & 12.33 & 0.1084 & 0.4504 & 8/16 \\
\bottomrule
\end{tabular}
\caption{Leaderboard B (COMET-Kiwi-22) --- end-to-end systems.}
\label{tab:comet-e2e}
\end{table}

\begin{table}[t]
\centering
\small
\begin{tabular}{@{}llRRRc@{}}
\toprule
\textbf{System} & \textbf{LTB-100 [95\% CI]} & \textbf{COMET-Kiwi} &
\textbf{IoU} & \textbf{$\tau$} & \textbf{Coverage} \\
\midrule
nllb-text-oracle v0.1.0    & 86.51 [85.3, 87.7] & 72.15 & 1.0000 & 1.0000 & 16/16 \\
deepl-text-oracle v0.1.0   & 84.89 [81.4, 88.0] & 69.77 & 1.0000 & 1.0000 & 6/16 \\
opus-mt-text-oracle v0.1.0 & 79.35 [77.9, 80.8] & 57.36 & 1.0000 & 1.0000 & 16/16 \\
\bottomrule
\end{tabular}
\caption{Leaderboard B (COMET-Kiwi-22) --- oracle-layout reference. Note the
ranking flip relative to Leaderboard A: NLLB-200 overtakes DeepL on COMET-Kiwi.}
\label{tab:comet-oracle}
\end{table}

\subsection{Key findings}

\paragraph{The chrF--COMET gap is largest on non-Latin pairs.} Rankings differ
between the two metrics. On the oracle-layout tables, DeepL leads on chrF
(78.20 vs.\ NLLB 73.71) but NLLB leads on COMET-Kiwi (86.51 vs.\ DeepL 84.89).
At the per-pair level the gap is concentrated on non-Latin scripts: for NLLB-200,
COMET-Kiwi exceeds chrF by roughly 31 points on en-kk (COMET 83.5 vs.\ chrF
52.2) and on en-th (COMET 76.0 vs.\ chrF 45.2), and by roughly 41 points on
en-ko (COMET 77.8 vs.\ chrF 36.9). chrF, a surface character-overlap metric,
systematically under-credits adequate translations into scripts and morphologies
unlike English; a learned metric does not. This divergence is exactly the
signal the dual-metric design is meant to expose---a single text metric would
have hidden it.

\paragraph{en$\rightarrow$uz is a structural zero.} Both open models effectively
fail Uzbek (Latin script): NLLB-200 scores chrF 0.07 and opus-mt scores chrF
0.04, with COMET-Kiwi around 0.31. The failure is consistent across metrics and
models, which indicates a genuine capability gap---Uzbek-Latin is not adequately
supported at the 600M-parameter open-model scale---rather than a metric
artifact. (A residual reference-quality concern on this pair cannot be fully
excluded and is noted in the methodology record.) A benchmark that surfaces such
a structural zero is doing useful work that a macro-average alone would obscure.

\paragraph{Open models diverge sharply by region.} A head-to-head between
opus-mt and NLLB-200 shows opus-mt is competitive with NLLB on European pairs
(within $\sim$3 COMET-Kiwi points on en-es, en-de, en-fr, en-ar, en-zh, en-ru,
with opus-mt slightly ahead on en-de and en-ru) but dramatically weaker on
several Asian and Central Asian pairs, where opus-mt's per-language models fall
tens of COMET-Kiwi points behind NLLB. Because opus-mt is Apache-2.0 / CC-BY-4.0
and NLLB-200 is CC-BY-NC-4.0, this per-pair profile---not a single macro
average---is the information a downstream practitioner actually needs.

\paragraph{Oracle-layout strongly outscores end-to-end.} The gap between the two
system classes is large. The best oracle-layout system reaches LTB-100 78.20
(chrF) by construction---it is handed perfect bounding boxes---while the two
non-trivial end-to-end systems measured reach far less: Florence-NLLB scores
39.79 (chrF) and 48.11 (COMET-Kiwi), and the zero-shot Qwen3-VL-2B-Instruct
14.82 and 18.43. The gap is driven by layout: end-to-end layout IoU is 0.1541
for Florence-NLLB and 0.1084 for Qwen3-VL, far below the oracle's 1.0000
ceiling. The identity baseline at
LTB-100~$\approx 50$ illustrates the composite's structure---it earns full IoU
and $\tau$ credit (it copies source boxes) and near-zero chrF (it does not
translate), so half of LTB-100's weight is unrealizable for it. End-to-end
document translation with faithful layout is, by this measurement, an open
problem; the headline number for a real system is the end-to-end number, not the
oracle-layout one.

\subsection{Florence-NLLB: an end-to-end OCR+MT pipeline}
\label{sec:florence}

Florence-NLLB is the first end-to-end OCR+MT pipeline runner in LTB. It is a
two-stage system: stage~1 uses Florence-2~\citep{xiao2024florence2} grounded OCR
(the \texttt{<OCR\_WITH\_REGION>} task) to detect text regions and their
bounding boxes directly from the source page image; stage~2 translates each
detected region with NLLB-200-distilled-600M~\citep{nllb2022}. Crucially the
bounding boxes come from Florence-2's predictions, \emph{not} from the
ground-truth annotation---so Florence-NLLB is a genuine end-to-end runner rather
than an oracle-layout one, and its layout score reflects a real region-detection
step.

Evaluated on the core-8 pairs (80 documents), Florence-NLLB scores LTB-100
39.79 [36.6, 42.9] on the chrF leaderboard and 48.11 [44.1, 51.7] on
COMET-Kiwi, ranking second of three end-to-end systems on both boards---behind
only the trivial identity baseline (50.38 / 50.47) and well ahead of the
zero-shot Qwen3-VL-2B-Instruct (14.82 / 18.43). Of its three components, layout
fidelity is the weakest: IoU is 0.1541---higher than Qwen3-VL's 0.1084, which is
consistent with Florence-2 being trained for document OCR grounding, but still
far below the oracle ceiling of 1.0. Florence-2 detects 5--13 regions per page
(mean 7.9, zero parser-failure placeholders), yet those regions only partially
overlap the annotated blocks, and unmatched ground-truth regions score zero.
Reading order is preserved reasonably well ($\tau = 0.81$), and text quality is
moderate---chrF 37.93, COMET-Kiwi 54.56---with the two metrics diverging in the
now-familiar pattern: the reference-free metric credits NLLB's fluent output
more generously than character overlap does, most visibly on the CJK pairs where
Florence-2's OCR is weakest (chrF 27.2 on en-zh and 29.4 on en-ja, against
38--47 on the Latin-script pairs). Florence-NLLB thus illustrates the paper's
central point concretely: a competent translation model paired with a
document-trained OCR model still leaves a wide gap to the oracle ceiling, and
that gap is dominated by imperfect layout extraction.

\section{Limitations}
\label{sec:limitations}

LTB is an evolving research artifact, and its methodology limitations are
tracked publicly in a methodology roadmap. We restate the substantive ones here.

\begin{itemize}
  \item \textbf{Single reference per region.} Each region currently has one
  reference translation. FLORES-200 references raise the quality bar to
  certified-translator level, but single-reference scoring still penalizes
  legitimate paraphrase variation. Multi-reference scoring (two references per
  document on the core 8) is deferred to v0.2 and requires a translation budget.
  \item \textbf{No human-evaluation correlation collected yet.} The
  infrastructure to correlate automatic LTB-100 with human Direct Assessment
  judgments is implemented (export of rating prompts and a correlation
  command), but no judgments have been collected. Until they are, the claim that
  LTB-100 tracks human judgment rests on the component metrics' established
  correlations~\citep{popovic2015chrf,rei2022cometkiwi}, not on LTB-specific
  evidence. Collecting $\sim$1000 judgments across systems is v0.2 scope and
  requires paid raters.
  \item \textbf{Dataset size.} 59 documents is small. Six core pairs are at
  $N=20$, and the extension pairs at $N=10$--$13$ have deliberately wide
  confidence intervals; they are coverage proofs, not ranking-grade samples.
  v0.2 targets 200 documents.
  \item \textbf{Category imbalance.} \texttt{ocr-document} is over-represented
  at 39 of 59 documents (66\%), so the headline number is weighted toward
  scanned real-world documents over forms, certificates, scientific papers, and
  contracts. v0.2 rebalances to $\geq 15$ documents per category.
  \item \textbf{Partial end-to-end coverage.} The end-to-end leaderboard covers
  the core-8 pairs only: both Florence-NLLB and Qwen3-VL-2B-Instruct are scored
  on 8 of 16 pairs. The oracle-layout tables are more complete but, by
  definition, do not measure layout extraction.
  \item \textbf{chrF as a text metric.} chrF is paraphrase- and adequacy-blind;
  the language-detection gate removes one specific artifact (Latin
  bleed-through) but not these deeper weaknesses. COMET-Kiwi-22 is provided as a
  parallel metric for this reason, and v0.2 will make it the leaderboard
  default.
  \item \textbf{Mined-reference quality.} The ML-curated subset required a
  script-validation pass to remove mislabeled references; residual
  reference-quality issues on individual pairs (notably en-uz) cannot be fully
  excluded.
  \item \textbf{No visual-fidelity or OCR round-trip metric yet.} v0.1 does not
  score destroyed figures, charts, or stamps, nor does it verify that text
  claimed in a submission was actually rendered. Both are v0.2 metrics
  (Section~\ref{sec:conclusion}).
\end{itemize}

\section{Conclusion and Future Work}
\label{sec:conclusion}

LayoutTranslateBench is, to our knowledge, the first public benchmark that
scores document translation, layout fidelity, and reading-order preservation
jointly under a single composite, LTB-100, with a reference-free second metric
on a parallel leaderboard. Version 0.1.7 provides 59 annotated documents across
9 categories and 16 language pairs. Our reference evaluation shows that the
chrF--COMET gap is concentrated on non-Latin scripts (validating the dual-metric
design), that en$\rightarrow$uz is a structural zero for 600M-class open models,
and that a large gap separates oracle-layout text-quality upper bounds from
realistic end-to-end systems---meaning faithful end-to-end document translation
remains an open problem.

The v0.2 roadmap is concrete:
\begin{itemize}
  \item \textbf{LPIPS visual-fidelity metric}~\citep{zhang2018lpips} --- a
  perceptual-similarity score on the non-text regions of the rendered output,
  catching destroyed figures, charts, stamps, and backgrounds; added to the
  composite at 10\% weight (Eq.~\ref{eq:ltb100v02}).
  \item \textbf{OCR round-trip metric} --- OCR the rendered output and compare
  the recognized text to the system's own declared translation, catching text
  claimed in a submission but not actually rendered; added at 10\% weight.
  \item \textbf{Multi-reference scoring} --- a second certified-translator
  reference per document on the core 8, reducing paraphrase penalization.
  \item \textbf{Human evaluation} --- Direct Assessment judgments across systems
  to empirically validate LTB-100's correlation with human judgment.
  \item \textbf{Dataset scale-up} --- to 200 documents, rebalanced to $\geq 15$
  per category, with broader end-to-end system coverage.
\end{itemize}
Submitted v0.1 results remain valid under the v0.2 composite and are displayed
with an explicit version label. The benchmark, dataset, leaderboard, and the
public methodology roadmap are released so that the community can both build
against LTB and argue with it.

\paragraph{Availability.} Code (Apache-2.0), dataset (CC-BY-4.0 /
CC-BY-SA-4.0), and the live leaderboard are available at
\url{https://github.com/Lawrenzho-bit/LayoutTranslateBench}; the benchmark
specification and dataset version are citable as~\citet{ltbench2026}.

% TODO: add Acknowledgments section with translator / contributor credits before arXiv submission

\bibliographystyle{plainnat}
\bibliography{refs}

\end{document}
